\hypertarget{ux5177ux8eabux667aux80fdux7684ux5c55ux671b}{%
\chapter{具身智能的展望}\label{ux5177ux8eabux667aux80fdux7684ux5c55ux671b}}

\hypertarget{ux5177ux8eabux667aux80fdux53d1ux5c55ux7684ux5173ux952eux8282ux70b9ux5c55ux671b}{%
\section{具身智能发展的关键节点展望}\label{ux5177ux8eabux667aux80fdux53d1ux5c55ux7684ux5173ux952eux8282ux70b9ux5c55ux671b}}

大语言模型诞生后，AI在数字世界快速发展；具身智能则意味着AI向物理世界拓展，打通现实世界中感知、推理与行动的闭环，带来智能形态的重要变革。下面是对未来具身智能发展的一些关键节点的展望。

\hypertarget{ux4e00ux5177ux8eabux667aux80fdux53d1ux5c55ux7684ux4e09ux4e2aux5173ux952eux8282ux70b9}{%
\subsection{一、具身智能发展的三个关键节点}\label{ux4e00ux5177ux8eabux667aux80fdux53d1ux5c55ux7684ux4e09ux4e2aux5173ux952eux8282ux70b9}}

\hypertarget{ux4e00chatgptux65f6ux523b}{%
\subsubsection{（一）ChatGPT时刻}\label{ux4e00chatgptux65f6ux523b}}

这个节点意味着具身智能具备泛化性，具备对物理世界的理解，在行动、操作等物理智能涉及的方面达到人类的水平，可以根据语言指令正确地完成物理世界中的特定元任务，如端菜品、叠衣服等。这里引入``物理图灵测试''的概念，即人类已经无法仅靠观察结果去判断完成工作的是人类还是机器人。

\hypertarget{ux4e8cux7269ux7406api}{%
\subsubsection{（二）物理API}\label{ux4e8cux7269ux7406api}}

在LLM的发展中，从Chatbot到Coding Agent是一次非常关键的形态跃迁，标志着模型能执行复杂的工程任务。在具身智能中，则意味着机器人可以通过物理API被调用，人或模型只需要描述要做的产品或任务是什么样的，具身智能就可以结合物理推理拆解任务、逐步执行动作等，直至做出对应产品或完成对应任务。这意味着工业、农业生产以及自然科学研究都可以实现闭环自动化。

\hypertarget{ux4e09ux7269ux7406ux81eaux6211ux6539ux8fdb}{%
\subsubsection{（三）物理自我改进}\label{ux4e09ux7269ux7406ux81eaux6211ux6539ux8fdb}}

人类与动物的物理智能的根本区别在于人类能发明和使用工具；LLM发展中，从单纯工程实现层面的Coding到自主研究、自我改进的AI for AI的转变也是一个关键的节点。对于具身智能而言，这个节点意味着具身智能不再只是在被告知产品具体形态下的生产执行者，而是``只有优化目标，自主探索做出什么样的东西实现''的研究者，即可以自主创造和使用工具、制造和改进具身智能自身，自主定义智能与物理世界之间的接口，将物理世界中的事物作为可闭环优化的对象，实现物理智能的自我改进和自主迭代。

\hypertarget{ux4e8cux4ecechatgptux65f6ux523bux5230ux7269ux7406apiux57faux672cux5f62ux6001ux4e0eux884cux4e1aux5e94ux7528}{%
\subsection{二、从ChatGPT时刻到物理API：基本形态与行业应用}\label{ux4e8cux4ecechatgptux65f6ux523bux5230ux7269ux7406apiux57faux672cux5f62ux6001ux4e0eux884cux4e1aux5e94ux7528}}

未来，我们希望基于具身智能的机器人能在现实任务里像人类一样完成工作，在各行各业被广泛应用。

\hypertarget{ux4e00ux65e5ux5e38ux751fux6d3bux5c42ux9762}{%
\subsubsection{（一）日常生活层面}\label{ux4e00ux65e5ux5e38ux751fux6d3bux5c42ux9762}}

具身智能将承担家务、照护、陪伴和家庭维修等任务。它不只是会执行固定动作的家电，而是能理解家庭成员的真实需求，感知复杂环境中的物品、老人、儿童和宠物状态，并在安全边界内自主规划行动，例如整理房间、准备餐食、照看老人、处理突发情况等。

\hypertarget{ux4e8cux670dux52a1ux4e1aux5c42ux9762}{%
\subsubsection{（二）服务业层面}\label{ux4e8cux670dux52a1ux4e1aux5c42ux9762}}

具身智能将进入餐饮、酒店、物流、医疗护理、零售和公共服务等场景。未来的机器人不仅能完成送餐、清洁、搬运、导览等重复性工作，也能根据顾客的语言、表情和行为做出灵活回应，从而让服务更稳定、更低成本，也让人类员工从高强度、低创造性的劳动中解放出来。

\hypertarget{ux4e09ux5de5ux4e1aux5c42ux9762}{%
\subsubsection{（三）工业层面}\label{ux4e09ux5de5ux4e1aux5c42ux9762}}

具身智能将实现工业生产的无人化。未来人类以自然语言等形式描述的需求可交由大模型得到可执行的生产方案，而实体生产过程可以完全由机器人自主完成。机器人可以操控流水线上的机器、监控生产状态，并对紧急状况及时作出反应。

\hypertarget{ux56dbux519cux4e1aux5c42ux9762}{%
\subsubsection{（四）农业层面}\label{ux56dbux519cux4e1aux5c42ux9762}}

具身智能将推动农业从依赖经验和人工劳动，走向精准化、自动化和数据化。机器人可以完成播种、除草、施肥、采摘、病虫害识别、土壤监测和农作物分拣等任务，并根据天气、土壤、水分和作物生长状态实时调整策略，从而提高产量、减少浪费，也缓解了农业劳动力不足的问题。

\hypertarget{ux4e09ux5177ux8eabux667aux80fdux7684ux81eaux6211ux6539ux8fdbux5177ux8eabux667aux80fdux63a8ux52a8ux667aux80fdux81eaux8eabux7684ux53d1ux5c55}{%
\subsection{三、具身智能的自我改进：具身智能推动智能自身的发展}\label{ux4e09ux5177ux8eabux667aux80fdux7684ux81eaux6211ux6539ux8fdbux5177ux8eabux667aux80fdux63a8ux52a8ux667aux80fdux81eaux8eabux7684ux53d1ux5c55}}

有认知理论认为，人类智能的形成并不是单纯依靠抽象符号推理完成的，而是在长期与环境互动的过程中逐渐发展出来的。人类婴儿通过看、听以及主动与世界互动，逐步建立对物体、空间、因果关系和他人行为的理解。如果这一理论是对的，那么实现具备人类全部能力的智能就需要具备具身的``感知---推理---行动''闭环，不仅要能观察世界、理解世界，还要能主动对世界采取行动并从中自主学习。具身智能可以在数据、模型、算力等多个维度对智能进行改进，进而实现物理世界中智能的递归自我改进。

\hypertarget{ux4e00ux5177ux8eabux667aux80fdux81eaux6211ux6539ux8fdbux7684ux65b9ux5f0fux63a8ux52a8ux5404ux7c7bux8981ux7d20ux6269ux5c55ux5b9eux73b0ux81eaux4e3bux5b66ux4e60}{%
\subsubsection{（一）具身智能自我改进的方式：推动各类要素扩展，实现自主学习}\label{ux4e00ux5177ux8eabux667aux80fdux81eaux6211ux6539ux8fdbux7684ux65b9ux5f0fux63a8ux52a8ux5404ux7c7bux8981ux7d20ux6269ux5c55ux5b9eux73b0ux81eaux4e3bux5b66ux4e60}}

\begin{enumerate}
\def\labelenumi{\arabic{enumi}.}
\tightlist
\item
  加快物理世界的迭代速度，支持算力不断扩张
\end{enumerate}

智能的发展依赖于算力的不断扩张，但这目前仍受制于人类的生产能力。具身智能可以通过闭环加快物理世界的迭代速度，通过扩大规模、优化流程等方式让可用算力持续扩张，避免智能发展被算力限制。

\begin{enumerate}
\def\labelenumi{\arabic{enumi}.}
\setcounter{enumi}{1}
\tightlist
\item
  突破数据局限，主动获取数据
\end{enumerate}

人类互联网的数据在规模、模态、内容范围上都是受限的，依赖人类把数据放进电脑里、``喂''给模型可能终究是不可持续的。当具身智能达到可以自主与环境互动的水平，就可以与现实世界交互，在交互中主动获取所需要的各种类型的环境数据（如在日常生活环境中主要是视觉、听觉、触觉和传感器数据等）以及验证自己的判断，随着具身智能的发展，其行动范围将不断扩大，可获得数据的分布的广度也不断增大，理论上这样可以获得可无限扩展的数据。模型主动行动以获取自己缺失的数据，正是自主学习的重要一环。

\begin{enumerate}
\def\labelenumi{\arabic{enumi}.}
\setcounter{enumi}{2}
\tightlist
\item
  扩展验证范围，构成物理世界中的优化闭环
\end{enumerate}

LLM在数学、代码和Agent任务上的提升主要依赖于基于可验证奖励的强化学习。具身智能正可以将可验证的范围拓宽到现实世界，从而让需要与现实世界交互的任务也可以构成可被模型处理的闭环，从而可以进行强化学习的``信用分配''和优化。

\begin{enumerate}
\def\labelenumi{\arabic{enumi}.}
\setcounter{enumi}{3}
\tightlist
\item
  实现自主学习的闭环
\end{enumerate}

自主学习即智能体主动在与环境的交互中获取数据和反馈，并以此主动更新自身认知。在具身智能或未来与大模型融合形成的下一代模型中，前者不仅包括用于深度学习的环境状态转移、动作等数据，还可以包括用于强化学习的验证信号；后者未来则可以结合利用给定数据或奖励优化模型的方法（如AI for AI等）。二者可以实现自主学习的闭环，算力的扩张则作为根本支撑，这也正好和《The Bitter Lesson》中提到的搜索与学习呼应。未来，对于具备具身能力的基础模型，或许我们可以在训练阶段用类似于RL的方法（如在需要探索环境的任务上训练）或其他方法训练出推理时可以在新环境中主动学习并快速适应，进而在其中出色地执行新任务的模型。

\hypertarget{ux4e8cux5177ux8eabux667aux80fdux81eaux6211ux6539ux8fdbux7684ux8868ux73b0ux591aux65b9ux9762ux6539ux8fdbux591aux73afux5883ux6269ux5c55}{%
\subsubsection{（二）具身智能自我改进的表现：多方面改进，多环境扩展}\label{ux4e8cux5177ux8eabux667aux80fdux81eaux6211ux6539ux8fdbux7684ux8868ux73b0ux591aux65b9ux9762ux6539ux8fdbux591aux73afux5883ux6269ux5c55}}

\begin{enumerate}
\def\labelenumi{\arabic{enumi}.}
\tightlist
\item
  实现多方面的自我改进
\end{enumerate}

LLM已经呈现出参与自我改进的迹象，如模型训练和智能体框架等方面，但其互动环境限于计算机内部。具身智能自我改进的涉及面则更加广泛，包括模型训练、物理结构设计、生产制造、测试与迭代等。具身智能不仅可以自主获取数据，还可以通过改进传感器和行动模块等方式，拓宽自身从环境获取数据、与环境交互的渠道，扩大感知和行动的空间。从这个角度看，后面将提到的发明实验技术或测量工具，本质上也可以视为具身智能自我改进的一环。

\begin{enumerate}
\def\labelenumi{\arabic{enumi}.}
\setcounter{enumi}{1}
\tightlist
\item
  将智能扩展到更多新环境
\end{enumerate}

当具身智能实现自主学习和自我改进，智能或许可以被扩展到更多新的环境。自主学习的智能体可以在新环境中不断探索、获取数据、更新认知，结合对新环境的认知和具身智能的自我改进能力，打造出在新环境中能力更强的具身智能，如纳米机器人、太空探索机器人、深海作业机器人、灾害救援机器人、实验室机器人以及医疗手术机器人等。它们可以在复杂的环境中实时感知、推理和行动，而非仅执行预编程路径。这些机器人自身也可以通过收集来自各自环境的交互数据、训练对应环境的模型、制造对应的硬件、改进获取数据的方式等，迭代得到对应环境更强的具身智能。在这个飞轮中，具身智能的``感知---推理---行动''闭环将被扩展到不同尺度、不同环境和不同任务中。

\hypertarget{ux53c2ux8003ux6587ux732e-165}{%
\subsection{参考文献}\label{ux53c2ux8003ux6587ux732e-165}}

\vspace{0.25em}
\begin{itemize}
\tightlist
\item
  Sutton, R. S. (2019). \href{http://www.incompleteideas.net/IncIdeas/BitterLesson.html}{The Bitter Lesson}. Incomplete Ideas.
\item
  Ha, D., \& Schmidhuber, J. (2018). \href{https://arxiv.org/abs/1803.10122}{World Models}. arXiv:1803.10122.
\item
  Physical Intelligence. (2026). \href{https://arxiv.org/abs/2604.15483}{π0.7: A Steerable Generalist Robotic Foundation Model with Emergent Capabilities}. arXiv:2604.15483.
\item
  Fan, J. (2026). \href{https://www.youtube.com/watch?v=3Y8aq_ofEVs}{Robotics' End Game: Nvidia's Jim Fan}. Sequoia Capital AI Ascent 2026 {[}Video{]}.
\end{itemize}

\clearpage

\hypertarget{ux5177ux8eabux667aux80fdux52a0ux901fux79d1ux5b66ux53d1ux73b0}{%
\section{具身智能加速科学发现}\label{ux5177ux8eabux667aux80fdux52a0ux901fux79d1ux5b66ux53d1ux73b0}}

\hypertarget{ux4e00ux7269ux7406apiux4e0bux7684ux79d1ux5b66ux53d1ux73b0}{%
\subsection{一、物理API下的科学发现}\label{ux4e00ux7269ux7406apiux4e0bux7684ux79d1ux5b66ux53d1ux73b0}}

对于自然科学和医学等领域，AI for Science的核心瓶颈在于数据与实验验证。当有潜力的方法和想法可以被大量提出，实验验证就成为了瓶颈；即便可以通过构建模拟来减少对湿实验的需求，构建模拟的难点依然不在于模型，而在于大量、高质量、可复现的扰动数据。只有具身智能才能自主地进行实验验证和获取数据。

\begin{enumerate}
\def\labelenumi{\arabic{enumi}.}
\tightlist
\item
  实现探索性研究的干湿闭环，解放人类研究者，降低探索成本
\end{enumerate}

当下，基于硬编码机器人的自动化实验室适配场景范围窄、定制耗时长、成本高昂，故往往出现在最不新颖的方向，难以加速长尾区域的科学发现。具身智能能像人一样灵活地执行实验，根据样品状态等感知信息实时调整行动；也能积累那些内蕴于``感觉''而不出现在文献中的隐性知识，如怎么构建厚度恰到好处的水膜等。细胞多组学数据获取中，当前测序和成像等已自动化，但具身智能还能处理组织解离、单细胞悬液制备、原代细胞处理、iPSC/类器官培养等需根据观测情况实时判断的预处理步骤；基因编辑的核心步骤已自动化，但具身智能还能打通下游克隆扩增、表型确认的链条。此外，软组织切片、动物实验、仪器间衔接、执行新实验流程、故障排错等都将由具身智能完成。具身智能落地后，AI Scientist将像从专用处理器向通用处理器转变那样，走向``通用科学家闭环''，探索性研究将不再大量依赖人力，从而研究人员能被解放出来进行更多创造性思考，具身智能的通用性也让单次探索的成本大大降低且可以闭环优化，有望带来更多新的科学发现。

\begin{enumerate}
\def\labelenumi{\arabic{enumi}.}
\setcounter{enumi}{1}
\tightlist
\item
  加快非标准化实验的速度
\end{enumerate}

当下，很多实验的通量已经增长了几个数量级。但由Amdahl定律，最终起决定性作用的依然是最慢的非标准化实验。人类往往每天高效工作的时间只有几个小时，且实验失败率较高；具身智能则可以7*24小时运行，且可以不断降低实验的失败率，从而加快单个实验的速度。

\begin{enumerate}
\def\labelenumi{\arabic{enumi}.}
\setcounter{enumi}{2}
\tightlist
\item
  绕过人类注意力瓶颈，扩大并行规模
\end{enumerate}

计算机科学告诉我们，若单个任务时间较长，则并行能让效率更高。具身智能可以让非标准化实验也被大规模并行，且不需要在遇到堵头、析出、污染、细胞状态不对等设计者没枚举的物理事件时就停下来（如果这些都需要人处理，那可达的并行度就被人的串行效率限制了）。除了少数存在因果关联或资源限制、需要串行的实验外，大部分实验可以实现在物理世界中的大规模并行。也有一些模型训练需要通过这种并行获取数据，如虚拟细胞模型中的不同类型扰动数据，这是人类分散式生产的数据给不出的。

\begin{enumerate}
\def\labelenumi{\arabic{enumi}.}
\setcounter{enumi}{3}
\tightlist
\item
  提高数据质量，让复现变易、模拟可行、模型更准
\end{enumerate}

专用模型（乃至模拟）依赖于大规模、高质量的数据，特别是因果扰动数据。而很多实验在不同人做的时候结果相差明显，如平板划线法中每个人划的圈数、方式等不同将直接影响菌落密度，这样带来的不一致的数据很难训练出好的模型。具身智能可实现对软组织、动物、微生物等的操作可溯源、可统一，让湿实验可被复现（复现的边界成本趋于零），减小批次效应的影响，还能收集那些在人类研究中常缺失或被忽略的负样本，为专用模型的训练提供高质量的数据。这将带来复利效应：更高质量的数据带来更好的模型，更好的模型预测更准确的结果，更准确的结果上收集更高质量的数据，从而极大加速研究。

\begin{enumerate}
\def\labelenumi{\arabic{enumi}.}
\setcounter{enumi}{4}
\tightlist
\item
  捕捉``意外''中的科学发现
\end{enumerate}

很多自然科学发现来源于``意外''。依赖硬编码机器人的实验室会忽略意外情况或停止，必须等待人类处理，瞬态现象可能很快就消失了；具身智能则可以主动感知、主动获取相关数据，比如换一个角度观察和获取数据等，这可能会催生出新的科学发现。

\hypertarget{ux4e8cux7269ux7406ux81eaux6211ux6539ux8fdbux4e0bux7684ux79d1ux5b66ux53d1ux73b0}{%
\subsection{二、物理自我改进下的科学发现}\label{ux4e8cux7269ux7406ux81eaux6211ux6539ux8fdbux4e0bux7684ux79d1ux5b66ux53d1ux73b0}}

当具身智能可以自我改进，则意味着它实现了像人类一样开始发明工具的跃迁，它不仅可以做实验，而且能够改造自己的实验接口、创造新的测量方法和实验技术，从而不断扩大从现实世界获取信息的能力，改造具身智能自身进行科学研究的方式和能力。

\begin{enumerate}
\def\labelenumi{\arabic{enumi}.}
\tightlist
\item
  改进实验操作和设备
\end{enumerate}

具备自我改进能力的具身智能可以自主改进实验的具体操作，如实时观察微小的形态学变化，结合力学反馈精准控制每个时刻换液的速度、抓取的力度，乃至进行局部微手术等，通过微调自身模型等形式实现自适应，大幅减少实验失败重做次数。具身智能还可以改进实验设备，设计更高效率的设备，提高实验的通量和准确性等。

\begin{enumerate}
\def\labelenumi{\arabic{enumi}.}
\setcounter{enumi}{1}
\tightlist
\item
  在新实验方法和测量工具的探索中，加快想法验证和工程迭代的速度
\end{enumerate}

实验方法和测量工具上的突破往往能带来重要的科学发现，因为它们能让我们绕过研究对象本身固有的复杂性，让我们直接看到问题的本质。尽管未来模型可能可以批量提出关于新实验方法和测量工具的想法（work和不work的都有），但验证和工程改进会成为瓶颈，无法构成闭环。具身智能可以以显著超过人类的速度在更大的空间并行试错和迭代改进，处理硬编码机器人无法处理的流程，优化有价值的方法，快速排除效果不佳的方法，技术从想法到在真实实验环境中可用的时间极大缩短、性能显著提高。

\begin{enumerate}
\def\labelenumi{\arabic{enumi}.}
\setcounter{enumi}{2}
\tightlist
\item
  探索新实验路径后，设计制造对应窄域自动化设备
\end{enumerate}

具身智能和窄域自动化并非相互独立的路径。在通用具身智能具备自我改进的能力后，其将有能力设计、制造和改进窄域自动化设备。这意味着在未来实验室的分工中，具身智能不总是自己做实验，而是自己负责探索性实验，找到稳定工作流后``编译''为高通量的自动化基础设施执行后续实验，自身则继续做其他探索并适当监测。

\begin{enumerate}
\def\labelenumi{\arabic{enumi}.}
\setcounter{enumi}{3}
\tightlist
\item
  基于物理交互中积累的直觉，推动新技术发明甚至新理论提出
\end{enumerate}

具身智能既有着对物理世界变化规律和空间结构的认知，又可以在大量实验中积累隐性知识或直觉。对于实验科学而言，其理解实验的难点和不会写入文献的隐性知识，解决了实验科学家和工程师的沟通难题，``知道什么值得做''的创造性科学直觉也可能与浸泡在实验室中的具身经验和隐性知识有关。而即便是对于理论物理等学科，人类也是在与物理世界的反复直接交互中，才从原始数据里抽象出了力、因果性等概念框架。大语言模型无法像牛顿那样将世界万物复杂的相互作用抽象为统一的``力''，因为它们对世界的认知来源于预训练数据，只能在已有数据包含的概念和框架中插值组合，这并不会得到全新的概念或理论，强化学习也不会让模型具备概念层次的创新能力。具身智能可以与现实世界主动交互，自主从物理世界获得已有框架之外的信息，这可能是实现概念级创新的必要条件。

\hypertarget{ux53c2ux8003ux6587ux732e-166}{%
\subsection{参考文献}\label{ux53c2ux8003ux6587ux732e-166}}

\vspace{0.25em}
\begin{itemize}
\tightlist
\item
  Szymanski, N. J., Rendy, B., Fei, Y., et al.~(2023). \href{https://www.nature.com/articles/s41586-023-06734-w}{An Autonomous Laboratory for the Accelerated Synthesis of Inorganic Materials}. Nature, 624, 86--91.
\item
  Lu, C., Lu, C., Lange, R. T., Foerster, J., Clune, J., \& Ha, D. (2024). \href{https://arxiv.org/abs/2408.06292}{The AI Scientist: Towards Fully Automated Open-Ended Scientific Discovery}. arXiv:2408.06292.
\item
  Canty, R. B., \& Abolhasani, M. (2026). \href{https://www.nature.com/articles/s41570-026-00847-2}{The Past, Present and Future of Self-Driving Laboratories}. Nature Reviews Chemistry, 10, 523--537.
\item
  Fan, J. (2026). \href{https://www.youtube.com/watch?v=3Y8aq_ofEVs}{Robotics' End Game: Nvidia's Jim Fan}. Sequoia Capital AI Ascent 2026 {[}Video{]}.
\end{itemize}

\clearpage
